\documentclass[11pt]{amsart}
\usepackage[margin=1.15in]{geometry}
\usepackage{amsmath,amssymb}
\usepackage[hidelinks]{hyperref}

\DeclareMathOperator{\ord}{ord}
\newcommand{\ellworH}{\ell^{\,\mathrm{wor-H}}}

\title{Literature-implied answer to the Hilbert \(L_2\) power-function problem}
\author{fa\_banach\_001 lane 5}
\date{2026-07-18}

\begin{document}
\maketitle

\section*{Status}

This is a lightweight literature-implied answer packet.  It records that
Open Problem 1 in arXiv:1011.3682 is answered affirmatively by the later
Hilbert \(L_2\) sampling theorem of Krieg--Ullrich, arXiv:1905.02516.

\section*{Source Question}

The source paper is Erich Novak and Henryk Wozniakowski, \emph{On the power
of function values for the approximation problem in various settings},
arXiv:1011.3682.

In Section 2.1, after Theorem 1, the paper asks:
\[
  \text{Suppose } r>\frac12. \text{ Is it true that }
  \ellworH(r,2)=1?
\]
Equivalently, in the source notation, whenever the \(n\)-th minimal worst-case
error using arbitrary linear information has polynomial rate \(r>1/2\) for a
Hilbert-space approximation problem measured in \(L_2\), do function values
have the same polynomial power?

\section*{Supporting Result}

The supporting paper is David Krieg and Mario Ullrich, \emph{Function values
are enough for \(L_2\)-approximation}, arXiv:1905.02516, published in
\emph{Foundations of Computational Mathematics} 21 (2021), 1141--1151.

Their main theorem gives constants \(C,c>0\) and \(k_n\ge c n/\log(n+1)\)
such that, for every reproducing kernel Hilbert space \(H\hookrightarrow L_2\),
\[
  e_n(H)^2 \le \frac{C}{k_n}\sum_{j\ge k_n} a_j(H)^2,
\]
where \(e_n\) are sampling numbers, i.e. errors using function values, and
\(a_n\) are approximation numbers, i.e. errors using arbitrary linear
information.  Their Corollary 1 states that if
\[
  a_n(H)\lesssim n^{-s}\log^\alpha(n)
  \qquad (s>1/2),
\]
then
\[
  e_n(H)\lesssim n^{-s}\log^{\alpha+s}(n),
\]
and in particular \(\ord(e_n)=\ord(a_n)\).

\section*{Implication for arXiv:1011.3682}

Fix \(r>1/2\).  Suppose the all-linear worst-case Hilbert \(L_2\) error has
source rate \(r\), i.e. \(\ord(a_n)=r\).  For every \(s<r\), the sequence
\(a_n\) is bounded by \(n^{-s}\) up to a constant, so Krieg--Ullrich's
corollary gives \(\ord(e_n)\ge s\).  Letting \(s\uparrow r\) yields
\(\ord(e_n)\ge r\).  Since function values are a restricted class of linear
information, \(a_n\le e_n\), hence \(\ord(e_n)\le \ord(a_n)=r\).  Therefore
\[
  \ord(e_n)=\ord(a_n)=r.
\]
In the power-function notation of arXiv:1011.3682, this is exactly
\[
  \ellworH(r,2)=1
  \qquad\text{for every }r>\frac12.
\]

\section*{Scope}

This packet answers only Open Problem 1 of the source paper.  It does not
settle the source's open problems for \(p\ne2\), Banach worst-case power
functions, randomized settings, average-case \(p\ne2\), or the fixed-\(n\)
supremum ratio problem.

\section*{Provenance}

The supporting paper explicitly answers the equivalent open question about
the equality of polynomial orders for Hilbert \(L_2\) sampling and arbitrary
linear information.  In the TeX source checked for this packet, it does not
appear to cite arXiv:1011.3682 directly; the connection to the source Open
Problem 1 is the power-function reformulation above.

\begin{thebibliography}{9}
\bibitem{NW}
E. Novak and H. Wozniakowski,
\emph{On the power of function values for the approximation problem in
various settings},
arXiv:1011.3682.

\bibitem{KU}
D. Krieg and M. Ullrich,
\emph{Function values are enough for \(L_2\)-approximation},
arXiv:1905.02516; Found. Comput. Math. 21 (2021), 1141--1151.
\end{thebibliography}

\end{document}
